	\subsection{Tools}
	
	The main technical tool we develop to prove Theorem \ref{thm:main0} is a structure theorem for set systems with small `doubling'. This result is similar in spirit to the famous Freiman-Ruzsa type theorems in additive combinatorics. These theorems describe the approximate structure of the subsets $A$ of a group with the property that the sum-set $A + A$ has size not much larger than $A$ (i.e. $A$ has small doubling). The theorem of Freiman and Ruzsa \cite{Freiman,Ruzsa} states that a subset of the integers with small doubling must be contained in a so-called generalized arithmetic progression of bounded rank (see also \cite{TV}). This classical result was subsequently generalized to abelian groups by Green and Ruzsa \cite{GR07} and to all groups by Breuillard, Green and Tao \cite{BGT12} (see also the survey \cite{BGT_Survey}). These results show that sets with small doubling must be close in structure to one of few natural examples. 
	
	Given a set-family $\mathcal{F}$, our measure for the size of $\mathcal{F}$ will be the dimension of the subspace $\langle \mathcal{F} \rangle$ spanned by the characteristic vectors of the sets in $\mathcal{F}$ over some field $\mathbb{F}$. 
	Let $\mathcal{F}\cdot\mathcal{F}=\{A\cap B:A,B\in\mathcal{F}\}$. Note that, by definition, $\mathcal{F} \subset \mathcal{F} \cdot \mathcal{F}$. 
	What can we say about $\mathcal{F}$ if the dimension of $\langle \mathcal{F} \cdot \mathcal{F} \rangle$ is not much larger than that of $\langle \mathcal{F} \rangle$? Observe that if $S$ is an atomic set-family, then $S \cdot S = S$.
	Our structure theorem shows that this is essentially the only possible example: any set-family $\mathcal{F}$ with small `doubling' must be close to being atomic. To make this more precise, we need the following definition. Given $i,j\in [n]$, say that $i$ and $j$ are {\em twins} for $\mathcal{F}$ if every $F \in \mathcal{F}$ either contains both $i,j$ or none of them, and there is at least one $F \in \mathcal{F}$ such that $i,j \in F$. Note that being twins is an equivalence relation (on the set of $i \in [n]$ which are contained in at least one $F \in \mathcal{F}$). A set of coordinates $T\subset [n]$ is called a \emph{set of twins} if any pair of elements in $T$ are twins, or $|T|=1$. Also, say that $T$ is a maximal set of twins if it is a complete equivalence class of the twins relation. We can now state our structure theorem. 
	
	
	\begin{theorem}\label{cor:stability}
		Let $\mathcal{F}\subset\{0,1\}^{n}$, let $\mathbb{F}$ be a field, and suppose that $\dim\langle \mathcal{F} \rangle = d$ and $\dim \langle \mathcal{F}\cdot \mathcal{F} \rangle=d+h$. Then $[n]$ can be partitioned into $d+1$ sets $A_1,\dots,A_d,B$ such that $A_{i}$ is a maximal set of twins for $\mathcal{F}$ for $i\in [d]$, and $\dim \langle \mathcal{F}|_{B} \rangle \leq 2h$.
	\end{theorem}
	
	\section{Small `doubling' and twins}
	
	In this section we establish Theorem \ref{cor:stability}. We will actually prove a more general statement about arbitrary vector spaces. Let us introduce some notation. 
	
	As usual, $\mathbb{Z}_{\ell}$ denotes the ring of integers modulo $\ell$, and if $p$ is a prime, we write $\mathbb{F}_p$ instead of $\mathbb{Z}_p$ to emphasize that it is also a field. Fix any commutative ring $\mathcal{R}$ with a unity  (in our case, $\mathcal{R}$ will be either a field or  $\mathbb{Z}_\ell$ for some positive integer $\ell$). For a vector $v \in \mathcal{R}^n$, we use $v(i)$ to denote the $i$th coordinate of $v$. The {\em support} of $v$ is $\{i \in [n] : v(i) \neq 0\}$. For $\mathcal{F}\subset \mathcal{R}^{n}$, we use $\langle \mathcal{F}\rangle$ to denote the \emph{span} (i.e. the set of all linear combinations) of the elements of $\mathcal{F}$. If $\mathcal{R}=\mathbb{Z}_{\ell}$, we might write $\langle \mathcal{F}\rangle_{\ell}$ instead of $\langle \mathcal{F}\rangle$ if $\mathcal{R}$ is not clear from the context. If $A\subset [n]$ and $v \in \mathcal{R}^n$, then $v|_{A}\in \mathcal{R}^{A}$ denotes the restriction of $v$ to the coordinates in $A$, and if $\mathcal{F}\subset \mathcal{R}^{n}$, then $\mathcal{F}|_{A}=\{v|_{A}: v \in \mathcal{F}\}$.
	
	Given vectors $v,w \in \mathcal{R}^n$, let $v\cdot w$ be the vector in $\mathcal{R}^n$ defined as $(v\cdot w)(i)=v(i)w(i)$ for $i\in [n]$. Note that if $v$ and $w$ are characteristic vectors of sets $A$ and $B$, then $v\cdot w$ is the characteristic vector of $A\cap B$. 
	For $V,W \subset \mathcal{R}^{n}$, let $V \cdot W = \{v \cdot w : v \in V, w \in W\}$. 
	Given $V \subset \mathcal{R}^{n}$ and $i,j\in [n]$, say that $i$ and $j$ are \emph{twins} for $V$ if  $v(i)=v(j)$ for all $v\in V$ and $v(i)\neq 0$ for at least one $v\in V$. Observe that if $V$ is a subspace generated by some family $\mathcal{F} \subset \{0,1\}^n$, then this definition of twins exactly coincides with the one given in the previous section.
	
	
	\begin{theorem}\label{lemma:stability}
		Let $\mathbb{F}$ be a field, $V<\mathbb{F}^{n}$, $d=\dim(V)$ and 
		$\dim(\langle V\cup (V\cdot V)\rangle)=d+h$. Then $[n]$ can be partitioned into $d+1$ sets $A_1,\dots,A_{d},B$ such that $A_{i}$ is a maximal set of twins for $V$ for each $i\in [d]$, and $\dim(V|_{B})\leq 2h$. 
	\end{theorem} 
	
	\begin{proof}
		For $i,j\in [n]$, say that $i$ and $j$ are \emph{siblings} for $V$ if there exists $\lambda\in\mathbb{F},\lambda\neq0$ such that $v(i)=\lambda v(j)$ for all $v\in V$, and $v(i)\neq 0$ for at least one $v\in V$. 
		Equivalently, $i$ and $j$ are siblings if for every basis $V_0$ of $V$ there exists $\lambda\in\mathbb{F},\lambda\neq0$ such that $v(i)=\lambda v(j)$ for all $v\in V_0$, and $v(i)\neq 0$ for at least one $v\in V_0$ (in which case $\lambda$ will be the same for all bases of $V$).  Note that if $\lambda = 1$, then $i$ and $j$ are twins.
		
		Let $v_1,\dots,v_d\in V$ be a basis of $V$, and let $M$ be the $d\times n$ matrix, whose rows are $v_1,\dots,v_d$. It is possible to choose the basis $v_1,\dots,v_d$ such that after possibly rearranging the columns of $M$, the restriction of $M$ to the first $d$ columns is a diagonal matrix. 
		
		We can assume without loss of generality that $M$ has no all-zero columns, i.e. that there is no coordinate $i \in [n]$ such that $v(i) = 0$ for all $v \in V$. Indeed, removing such an index does not change the dimension of $V$ and $\langle V\cup (V\cdot V)\rangle$, and we may include it in $B$ without increasing the dimension of $V|_{B}$.
		
		Note that the indices $i\in [d]$ and $j\in [n]\setminus [d]$ are siblings for $V$ if and only if $v_i(j)\neq 0$ and the $j$-th column of $M$ contains exactly one nonzero entry (namely, the entry $v_i(j)$). For $i\in [d]$, let $S_i$ contain all the siblings of $i$, and let $B'=[n]\setminus (S_1\cup\dots\cup S_d)$. Note that for every $j\in B'$, the $j$-th column of $M$ contains at least two nonzero entries.
		
		Let $r=\dim(V|_{B'})$. Choose some subset $C\subset B'$ such that the columns of $C$ form a basis of the vector space spanned by the columns of $B'$. Then $|C|=r$ and $\dim(V|_{C})=r$ (because the row space has the same dimension as the column space). See Figure \ref{fig1} for an illustration of the matrix $M$ and sets $S_1,\dots,S_d,B',C$. Furthermore, let $w_{i}=v_{i}|_{C}$ for $i\in [d]$, and observe that $w_1,\dots,w_d$ spans $\mathbb{F}^{C}$ (because $|C| = r$ and $\dim(V|_{C})=r$). For $c\in C$, let $\mathds{1}_c\in\mathbb{F}^{C}$ be the characteristic vector of the single element set $\{c\}$.
		
		\begin{claim}\label{claim:c'}
			For every $c\in C$ there exists $\alpha\in \mathbb{F}$, $c'\in C, c'\neq c$, and two vectors of coefficients $x,y\in \mathbb{F}^{d}$ such that $x$ and $y$ have disjoint supports and 
			$$\mathds{1}_{c}+\alpha\mathds{1}_{c'}=\left(\sum_{i=1}^d x(i)w_i\right)\cdot \left(\sum_{i=1}^d y(i)w_i\right).$$
		\end{claim}	
		\noindent
		Here, as always, $\cdot$ denotes the coordinate-wise product of vectors.
		
		\begin{figure}
			\centering
			\begin{tikzpicture}
				\node at (-3.5,0.7) {$v_1$};
				\node at (-3.5,0.233) {$v_2$};
				\node at (-3.5,-0.233) {$v_3$};
				\node at (-3.5,-0.7) {$v_4$};
				\matrix at (0,0) [matrix of math nodes,left delimiter=(,right delimiter=),row sep=0cm,column sep=0cm] (m) {
					1 &   &   &   &   & 2 &   &   &   & 2 &   & 1\\
					& 1 &   &   & 2 & 1 &   & 2 & 1 &   & 2 & 2\\
					&   & 1 &   &   &   &   & 1 & 1 &   &   & 1\\
					&   &   & 1 &   &   & 1 &   & 1 &   &   & 2\\};
				\draw[dashed] (-0.9,-1) -- (-0.9,1);
				\matrix (m) at (8,0) [matrix of math nodes,left delimiter=(,right delimiter=),row sep=0cm,column sep=0cm] (m) {
					1 & 2 &   &   &   &   &   &   & 2 &   &   & 1\\
					&   & 1 & 2 & 2 &   &   &   & 1 & 2 & 1 & 2\\
					&   &   &   &   & 1 &   &   &   & 1 & 1 & 1\\
					&   &   &   &   &   & 1 & 1 &   &   & 1 & 2\\};
				\draw[dashed] (8.9,-1) -- (8.9,1);
				\draw[dashed] (10.25,-1) -- (10.25,1);
				\draw [thick,decorate,decoration={calligraphic brace,amplitude=4pt}] (5.4,1) -- (6.2,1);
				\node at (5.8,1.5) {$S_1$};
				\draw [thick,decorate,decoration={calligraphic brace,amplitude=4pt}] (6.4,1) -- (7.45,1);
				\node at (6.925,1.5) {$S_2$};
				%
				\node at (7.8,1.5) {$S_3$};
				\draw [thick,decorate,decoration={calligraphic brace,amplitude=4pt}] (8.1,1) -- (8.8,1);
				\node at (8.45,1.5) {$S_4$};
				\draw [thick,decorate,decoration={calligraphic brace,amplitude=4pt}] (8.9,1) -- (10.25,1);
				\node at (9.425,1.5) {$C$};
				\draw [thick,decorate,decoration={calligraphic brace,amplitude=4pt}]  (10.6,-1) -- (8.9,-1);
				\node at (9.6,-1.5) {$B'$};
			\end{tikzpicture}
			\caption{On the left, a matrix $M$ over $\mathbb{F}_3$, where blank entries denote 0. On the right, we rearranged the columns of $M$ for better visualization.}
			\label{fig1}
		\end{figure}
		
		
	
		\begin{proof}
			Let $K\subset [d]$ be the set of indices $i \in [d]$ such that $w_{i}(c)\neq 0$. Recall that $|K|\geq 2$, as every column indexed by an element of $C\subset B'$ contains at least two nonzero entries. As $w_1,\dots,w_d$ span $\mathbb{F}^{C}$, we can choose an $r$-element subset of $\{w_1,\dots,w_d\}$ forming a basis of $\mathbb{F}^{C}$. Without loss of generality, suppose that $w_1,\dots,w_r$ is a basis. Then, for every $u\in C$, we can write $\mathds{1}_u=\sum_{i=1}^r\lambda_{u,i} w_i$ with suitable $\lambda_{u,1},\dots,\lambda_{u,r}\in\mathbb{F}$. Consider three cases.
			
			\begin{description}
				\item[Case 1.] There exists $k\in K$ such that $k\not\in [r]$. 
				
				Recall that, by definition of $K$, $w_k(c)\not =0$. Set $\alpha=0$, $x=(\lambda_{c,1},\dots,\lambda_{c,r},0,\dots,0) \in \mathbb{F}^d$, and define $y \in \mathbb{F}^d$ as 
				$$y(i)=\begin{cases}\frac{1}{w_k(c)} &\mbox{if }i=k,\\
					0                &\mbox{otherwise}.
				\end{cases}$$
				Then $\sum_{i=1}^d x(i)w_i=\mathds{1}_{c}$ and $\sum_{i=1}^d y(i)w_i=\frac{1}{w_k(c)}w_k$, which implies that $(\sum_{i=1}^d x(i)w_i)\cdot (\sum_{i=1}^d y(i)w_i)=\mathds{1}_c$. Furthermore, the support of $x$ is contained in $[r]$, while the support of $y$ is $\{k\}$, so $x$ and $y$ have disjoint supports. We see that $\alpha,x,y,c'$ satisfy the requirements for any choice of $c'\in C$, $c'\neq c$. 
				
				\item[Case 2.] $K\subset [r]$ and there exists $k\in K$ such that $\lambda_{c,k}=0$. 
				
				Define $\alpha,x,y$ the same way as in the previous case. Note that the only difference is that the support of $x$ is now contained in $[r]\setminus\{k\}$, so $x$ and $y$ still have disjoint supports. Therefore, $\alpha,x,y,c'$ suffice for any $c'\in C$, $c'\neq c$.
				
				\item[Case 3.] $K\subset [r]$ and $\lambda_{c,k}\neq 0$ for every $k\in K$.
				
				We claim that there exists some $c'\in C,c'\neq c$ such that not all $|K|$ coefficients $\lambda_{c',k}$ for $k\in K$ vanish. 
				Indeed, suppose by contradiction that $\lambda_{c',k} = 0$ for all $k\in K$ and $c' \in C \setminus \{c\}$. By the definition of the coefficients $\lambda_{c',i}$, this means that $\mathds{1}_{c'}$ belongs to the span of $(w_j)_{j \in [r] \setminus K}$ for every $c' \in C \setminus \{c\}$. But the $r-1$ vectors $(\mathds{1}_{c'})_{c'\in C \setminus \{c\}}$ are linearly independent, while $r - |K| \leq r-2$, so this is impossible. 
				%
				
				Choose $c' \in C \setminus \{c\}$ and $k\in K$ such that $\lambda_{c',k}\neq 0$. Let $\beta=-\lambda_{c,k}/\lambda_{c',k}$, 
				$$x=(\lambda_{c,1}+\beta\cdot \lambda_{c',1},\dots,\lambda_{c,r}+\beta\cdot\lambda_{c',r},0,\dots,0),$$ and define $y$ as
				$$y(i)=\begin{cases}\frac{1}{w_k(c)} &\mbox{if }i=k,\\
					0                &\mbox{otherwise}
				\end{cases}.$$
				Then $\sum_{i=1}^d x(i)w_i=\mathds{1}_c+\beta \mathds{1}_{c'}$ and $\sum_{i=1}^d y(i)w_i=\frac{1}{w_k(c)}w_k$. Therefore, we have 
				$$\left(\sum_{i=1}^d x(i)w_i\right)\cdot \left(\sum_{i=1}^d y(i)w_i\right)=\mathds{1}_c+\frac{\beta w_k(c')}{w_k(c)}\mathds{1}_{c'}.$$ Also, by the choice of $\beta$, we have $x(k)=0$, so the support of $x$ is contained in $[r]\setminus\{k\}$, while the support of $y$ is $\{k\}$. Therefore, $c',\alpha=\frac{\beta w_k(c')}{w_k(c)},x,y$ satisfy the requirements of the claim.
			\end{description}
			This completes the proof of Claim \ref{claim:c'}. See Figure \ref{fig2} for an illustration of $M|_{C}$ and the three cases.
		\end{proof}
		
		\begin{figure}
			\centering
			\begin{tikzpicture}
				\matrix at (0,0) [matrix of math nodes,left delimiter=(,right delimiter=),row sep=0cm,column sep=0cm] {
					& 1 & 1 \\
					1 & 1 & 2 \\
					2 &   & 1 \\
					2 &   &   \\
					1 &   &   \\
					2 &   &   \\};
				\draw[dashed] (-1,0) -- (1,0);
				\node at (-1.5,1.15) {$w_1$};
				\node at (-1.5,0.2) {$w_r$};
				\node at (-1.5,-0.2) {$w_{r+1}$};
				\node at (-1.5,-1.15) {$w_{d}$};
				\node at (-0.45,1.8) {$c_1$};
				\node at (0,1.8) {$c_2$};
				\node at (0.45,1.8) {$c_3$};
				\draw [thick,decorate,decoration={calligraphic brace,amplitude=4pt}] (0.9,-1.7) -- (-0.9,-1.7);
				\node at (0,-2) {$C$};
				
				\node[anchor=west] at (3,-0.6) {\large $\mathds{1}_{c_3}=w_1+2w_2+2w_3$};
				\node[anchor=west] at (3,0.6) {\large $\mathds{1}_{c_1}=w_1+2w_2+w_3$};
				\node[anchor=west] at (3,0) {\large $\mathds{1}_{c_2}=\ \ \ \ \ \ \ \ \ \ w_2+w_3$};
			\end{tikzpicture}
			\caption{An example of the matrix $M|_{C}$ in the field $\mathbb{F}_3$, where blank entries denote 0. The column $c=c_1$ falls into Case 1., $c_2$ falls into Case 2., $c_3$ falls into Case 3.}
			\label{fig2}
		\end{figure}
		
		For $c\in C$, let $x,y\in\mathbb{F}^{d}$ be the vectors of coordinates  satisfying the conditions of the previous claim. Also, let $X=\sum_{i=1}^d x(i)v_i$, $Y=\sum_{i=1}^d y(i)v_i$, and set $z_{c}=X\cdot Y$. Then $z_{c}\in V\cdot V$ and $z_{c}|_{C}=\mathds{1}_c+\alpha\mathds{1}_{c'}$ for some $c'\in C,c'\neq c$ and $\alpha\in \mathbb{F}$, as guaranteed by Claim \ref{claim:c'}. Observe that the support of $X$ on $[n]\setminus B'$ is $\bigcup_{i\in [d], x(i)\neq 0}S_i$, while the support of $Y$ on $[n]\setminus B'$ is $\bigcup_{i\in [d], y(i)\neq 0}S_i$ (indeed, recall that $[n]\setminus B'$ is the union of $S_1,\dots,S_d$, and every column in $S_i$ has a non-zero entry in the $i$th row and zeros in all other rows).
		Hence, as $x$ and $y$ have disjoint supports, $X$ and $Y$ have disjoint supports on $[n]\setminus B'$. This means that $z_c=X\cdot Y$ vanishes on $[n]\setminus B'$.
		
		 Let $W$ be the vector-space generated by the vectors $(z_c)_{c\in C}$, and let us record that $W$ vanishes on $[n]\setminus B'$.
		\begin{claim}
			$\dim(W)\geq r/2.$
		\end{claim}
		\begin{proof}
			Suppose that $\dim(W) < r/2$, and without loss of generality, let $D\subset C$ such that $\{z_c\}_{c\in D}$ is a basis of $W$ and $|D|<r/2 = |C|/2$. 
			Recall that for each $c \in C$, $z_{c}|_C = \mathds{1}_c+\alpha\mathds{1}_{c'}$ for some $c'\in C \setminus \{c\}$ and $\alpha\in \mathbb{F}$.
			As $z_{c}|_C$ vanishes on all but at most two coordinates, there exists $u\in C$ such that $z_c$ vanishes on $u$ for every $c\in D$. But $z_u$ does not vanish on $u$, so $z_u$ is not in the span of $\{z_c\}_{c\in D}$, contradicting that $\{z_c\}_{c\in D}$ is a basis for $W$.
		\end{proof}
		
		Let $I$ be the set of indices $i \in [d]$ such that the entries in $v_i|_{S_i}$ are {\em not} all equal, that is, $S_i$ is not a set of twins. For $i\in [d]$, let $A_{i}$ be a maximal set of twins in $S_{i}$, then $S_{i}\neq A_{i}$ if and only if $i\in I$. If $i\in I$, and $v_i|_{A_i}$ is the constant $s$ vector, then define $v_{i}'=sv_{i}-v_{i}\cdot v_{i}\in \langle V\cup (V\cdot V)\rangle$. Observe that $v_{i}'$ vanishes on $A_{i}$ and on $S_{j}$ for every $j\in [d]\setminus\{i\}$, and $v_{i}'$ does not vanish on $S_{i}\setminus A_{i}$. 
		Crucially, this implies that the $d+|I|$ vectors $v_1,\dots,v_{d},\{v_{i}'\}_{i\in I}$ are not only linearly independent, but their restrictions on $[n]\setminus B'$ are also linearly independent. 
		This means that setting $V'=\langle \{v_{i}:i\in [d]\}\cup\{v_{i}':i\in I\}\rangle$, we have $\dim(V')=d+|I|$, and the only element of $V'$ vanishing on $[n]\setminus B'$ is $\mathbf{0}$. On the other hand, every element of $W$ vanishes on $[n]\setminus B'$, which gives $W\cap V'=\{\mathbf{0}\}$.  See Figure \ref{fig3} for an illustration of the vectors $v_1,\dots,v_d,\{v'_i\}_{i\in I},\{z_c\}_{c\in C}$.
		
		But then as $V'+W < \langle V\cup (V\cdot V)\rangle$, we have
		$$\dim(\langle V\cup (V\cdot V)\rangle)\geq\dim (V'+W)=\dim(V')+\dim(W)\geq d+|I|+r/2.$$
		On the other hand, $\dim(\langle V\cup (V\cdot V)\rangle) = d+h$ by assumption.
		Therefore, $|I|+r/2\leq h$. Let $B=B'\cup\bigcup_{i\in I}(S_{i}\setminus A_{i})$, then $\dim(V|_{B})\leq r+|I|\leq 2h$, so the sets $A_{1},\dots,A_{d},B$ satisfy the desired properties.
	\end{proof}
	
	\begin{figure}
		\centering
		\begin{tikzpicture}
			\node at (4.6,1.8) {$v_1$};
			\node at (4.6,1.333) {$v_2$};
			\node at (4.6,0.866) {$v_3$};
			\node at (4.6,0.4) {$v_4$};
			\node at (4.6,-0.066) {$z_{c_1}$};
			\node at (4.6,-0.533) {$z_{c_2}$};
			\node at (4.6,-1) {$z_{c_3}$};
			\node at (4.6,-1.466) {$v_{1}'$};
			\node at (4.6,-1.933) {$v_2'$};
			\draw[rounded corners,color=white,fill=white!80!red] (5.3,2.3) rectangle (5.75,-2.3);
			\draw[rounded corners,color=white,fill=white!80!blue] (6.7,2.3) rectangle (7.5,-2.3);
			\draw[rounded corners,color=white,fill=white!80!green] (7.55,2.3) rectangle (8,-2.3);
			\draw[rounded corners,color=white,fill=white!80!brown] (8.02,2.3) rectangle (8.82,-2.3);
			\node at (5.525,-2.9) {$A_1$};
			\node at (7.18,-2.9) {$A_2$};
			\node at (7.81,-2.9) {$A_3$};
			\node at (8.42,-2.9) {$A_4$};
			\draw [thick,decorate,decoration={calligraphic brace,amplitude=4pt}] (7.5,-2.4) -- (6.7,-2.4);
			\draw [thick,decorate,decoration={calligraphic brace,amplitude=4pt}] (8.82,-2.4) -- (8.02,-2.4);
			
			\matrix at (8,0) [matrix of math nodes,left delimiter=(,right delimiter=),row sep=0cm,column sep=0cm] {
				1 & 2 &   &   &   &   &   &   & 2 &   &   & 1\\
				&   & 1 & 2 & 2 &   &   &   & 1 & 2 & 1 & 2\\
				&   &   &   &   & 1 &   &   &   & 1 & 1 & 1\\
				&   &   &   &   &   & 1 & 1 &   &   & 1 & 2\\
				&   &   &   &   &   &   &   & 1 &   &   & 1\\
				&   &   &   &   &   &   &   &   & 1 & 2 & 1\\
				&   &   &   &   &   &   &   &   & 2 & 1 & 2\\	  
				& 1 &   &   &   &   &   &   & 1 &   &   &  \\
				&   &  1 &   &  &   &   &   & 1 &   & 1 &  \\	  
			};
			
			\draw[dashed] (8.9,-2.3) -- (8.9,2.3);
			\draw[dashed] (10.25,-2.3) -- (10.25,2.3);
			\draw[dashed] (5.2,0.23) -- (10.8,0.23);
			\draw[dashed] (5.2,-1.3) -- (10.8,-1.3);
			\draw [thick,decorate,decoration={calligraphic brace,amplitude=4pt}] (5.3,2.4) -- (6.2,2.4);
			\node at (5.75,2.9) {$S_1$};
			\draw [thick,decorate,decoration={calligraphic brace,amplitude=4pt}] (6.4,2.4) -- (7.45,2.4);
			\node at (6.925,2.9) {$S_2$};
			\node at (7.81,2.9) {$S_3$};
			\draw [thick,decorate,decoration={calligraphic brace,amplitude=4pt}] (8.02,2.4) -- (8.82,2.4);
			\node at (8.42,2.9) {$S_4$};
			\node at (9.1,2.9) {$c_1$};
			\node at (9.525,2.9) {$c_2$};
			\node at (9.95,2.9) {$c_3$};
			\draw [thick,decorate,decoration={calligraphic brace,amplitude=4pt}]  (10.6,-2.4) -- (8.9,-2.4);
			\node at (9.75,-2.9) {$B'$};
			\draw [thick,decorate,decoration={calligraphic brace,amplitude=4pt}] (11.3,2.3) -- (11.3,0.2);
			\node at (12.8,1.25) {$M$};
		\end{tikzpicture}
		\caption{An example of a matrix $M$ over $\mathbb{F}_3$, together with the vectors $z_{c}$ for $c\in C=\{c_1,c_2,c_3\}$, and the vectors $v_{i}'$ for $i\in I$.}
		\label{fig3}
	\end{figure}
	
	
	
	\section{\texorpdfstring{$k$}{k}-closed families are atomic}
	
	In this section, we prove a theorem which implies Theorem \ref{thm:main0} after a small amount of work. Before stating it, we need some additional notation.  Say that $\mathcal{F}$ is \emph{non-reducible} if $\mathcal{F}$ does not vanish on any of the coordinates (namely, if there is no $i$ such that $v(i) = 0$ for all $v \in \mathcal{F}$). Recall that $\mathcal{F} \cdot \mathcal{F}$ is the set of all $v \cdot w$, $v,w \in \mathcal{F}$, where $v \cdot w$ is the coordinate-wise product. We also put $v^k = v \cdot \dotsc\cdot v$ and $\mathcal{F}^{k}=\mathcal{F}\cdot \dotsc \cdot \mathcal{F}$, where the products contain $k$ terms. Finally, let $||v||=\sum_{i=1}^{n}v(i)$. 
	
	We say that a set $\mathcal{F}\subset \mathbb{Z}^{n}_{\ell}$ is \emph{$k$-closed} if $||v||=0$ for every $1\leq i\leq k$ and $v\in \mathcal{F}^{i}$. Note that if $\mathcal{F} \subset \{0,1\}^n$, then this is the same as saying that the intersection of any $k$ not necessarily distinct sets from $\mathcal{F}$ is divisible by $\ell$.  Let us collect some simple properties of the coordinate-wise product and $k$-closedness.
	
	\begin{claim}\label{claim:k-closed}
		\begin{enumerate}
			
			\item  If $\mathcal{F},\mathcal{F}'\subset \mathbb{Z}_{\ell}^{n}$, then $\langle \mathcal{F}\rangle\cdot \langle \mathcal{F}'\rangle\subset \langle \mathcal{F}\cdot\mathcal{F}'\rangle$.
			
			\item 	If $\mathcal{F}\subset \mathbb{Z}_{\ell}^{n}$ is $k$-closed, then $\langle \mathcal{F}\rangle$ is also $k$-closed.
			
			\item If $\mathcal{F}$ is $k$-closed, then $\mathcal{F}\cdot\mathcal{F}$ is $\lfloor k/2\rfloor$-closed.
		\end{enumerate}
	\end{claim}
	
	\begin{proof}
		1. Let $v\in\langle\mathcal{F}\rangle$ and $w\in\langle\mathcal{F}'\rangle$. Then there exists $a_x\in\mathbb{Z}_{\ell}$ for every $x\in\mathcal{F}$ such that $v=\sum_{x\in\mathcal{F}}a_x x$, and similarly, there exists $b_y\in\mathbb{Z}_{\ell}$ for every $y\in\mathcal{F}'$  such that $w=\sum_{y\in\mathcal{F}'}b_y y$. But then
		$$v\cdot w=\left(\sum_{x\in \mathcal{F}}a_x x\right)\cdot \left(\sum_{y\in\mathcal{F}'}b_y y\right)=\sum_{x\in\mathcal{F}}\sum_{y\in\mathcal{F}'}a_xb_y(x\cdot y)\in \langle \mathcal{F}\cdot\mathcal{F}'\rangle.$$
		
		2. Let $i\in [k]$. By repeatedly applying 1., we can write $\langle\mathcal{F}\rangle^{i}\subset \langle \mathcal{F}^{i}\rangle$.
		As $\mathcal{F}$ is $k$-closed, we have $||v||=0$ for every $v\in\mathcal{F}^{i}$. 
		As $|| \cdot ||$ is a linear function,
%
		we have $||w||=0$ for every $w\in \langle \mathcal{F}^{i}\rangle$. Hence, $||w||=0$ for every $1\leq i\leq k$ and $w\in \langle \mathcal{F}\rangle^{i}$, meaning that $\langle \mathcal{F}\rangle$ is $k$-closed.
		
		3. This follows trivially from the definition.
	\end{proof}
	
	\noindent
	Furthermore, let us record some simple properties of the twin relation. 
	\begin{claim}\label{claim:twins}
		Let $\mathcal{F}\subset \{0,1\}^{n}$.
		\begin{enumerate}
			\item If $\mathcal{F}$ is non-reducible, then the maximal sets of twins for $\mathcal{F}$ form a partition of $[n]$.
			
			\item For every $k\geq 1$, the family $\bigcup_{i=1}^{k} \mathcal{F}^{i}$ has the same pairs of twins as $\mathcal{F}$.
			
		\end{enumerate}
	\end{claim}
	\begin{proof}
		1. This is a consequence of the fact that the twin relation is an equivalence relation.
		
		2. Let $a,b\in [n]$, $a\neq b$. If $a$ and $b$ are twins in $\mathcal{F}$, then $v(a)=v(b)$ for every $v\in\mathcal{F}$, which implies that $v_1(a)\dots v_i(a)=v_1(b)\dots v_i(b)$ for every $i\in [k]$ and $v_1,\dots,v_i\in\mathcal{F}$. Hence, $a$ and $b$ are twins in $\mathcal{F}^{i}$ as well, and so in $\bigcup_{i=1}^{k}\mathcal{F}^i$. Also, if $a$ and $b$ are not twins, then there exists some $v\in\mathcal{F}\subset \bigcup_{i=1}^{k}\mathcal{F}^i$ such that $v(a)\neq v(b)$, so $a$ and $b$ are not twins in $\bigcup_{i=1}^{k}\mathcal{F}^i$.
	\end{proof}
	Recall that $S(n,\ell) \subset 2^{[n]}$ is the atomic set-family with $\lfloor n/\ell \rfloor$ atoms of size $\ell$ each. 
	The main result of this section is the following variant of Theorem \ref{thm:main0}. We show that if $\mathcal{F}\subset 2^{[n]}$ is such that the intersection of any $k$ not necessarily distinct elements of $\mathcal{F}$ has size divisible by $\ell$, then $|\mathcal{F}|\leq 2^{\lfloor n/\ell\rfloor}$, given $k$ is sufficiently large with respect to $\ell$. We also show that if $\mathcal{F}$ is close to being extremal, then $\mathcal{F}$ must be a subfamily of (an isomorphic copy of) $S(n,\ell)$. 
	\begin{theorem}\label{thm:main}
		Let $\ell$ be a positive integer, then there exists $k$ such that the following holds. Let $\mathcal{F}\subset \{0,1\}^{n}$ such that $\mathcal{F}$ is $k$-closed over $\mathbb{Z}_{\ell}$. Then $|\mathcal{F}|\leq 2^{\lfloor n/\ell\rfloor}$. Also, if $|\mathcal{F}|> 2^{\lfloor n/\ell\rfloor-1}$, then $[n]$ can be partitioned into sets $A_{1},\dots,A_{d},A'$ such that $A_{i}$ is a maximal set of twins for $\mathcal{F}$ for $i\in [d]$, $|A_{i}|=\ell$, $|A'|\leq \ell-1$, and $\mathcal{F}$ vanishes on $A'$. 
	\end{theorem}
	We start preparing the proof of Theorem \ref{thm:main} with a sequence of lemmas. 
	The Eventown theorem mentioned in the introduction can be easily extended to vector spaces, where we replace intersections with scalar products. We will make use of the following simple extension of this in which we consider certain bilinear forms instead of scalar product.
	
	\begin{lemma}\label{lemma:bilinear}
		Let $\mathbb{F}$ be a field, let $b_1,\dots,b_d\in \mathbb{F}$, let $z$ be the number of zeros among $b_1,\dots,b_d$, and let $b:\mathbb{F}^{d}\times\mathbb{F}^{d}\rightarrow \mathbb{F}$ be the bilinear form defined as $b(v,w)=\sum_{i=1}^{d}b_i v(i)w(i)$. Let $V<\mathbb{F}^{d}$ such that $b(v,w)=0$ for every $v,w\in V$. Then $\dim(V)\leq \frac{1}{2}(d+z)$.
	\end{lemma}
	
	\begin{proof}
		Let $M$ be the $d\times d$ diagonal matrix with diagonal entries $b_{1},\dots,b_{d}$, and let $W=\{Mv:v\in V\}<\mathbb{F}^{d}$. Then $\dim(\mbox{ker }M)=z$, so $\dim(W)\geq\dim(V)-z$. 
		By definition, $V$ and $W$ are orthogonal spaces (with respect to the standard inner product). Therefore, $\dim(V)+\dim(W)\leq d$, which implies $\dim(V)\leq \frac{1}{2}(d+z)$.
	\end{proof}
	
	In what comes, we show that if $\ell=p^{\alpha}$ is a prime power, and  $\mathcal{F}\subset \{0,1\}^n$ is $k$-closed over $\mathbb{Z}_{\ell}$ for some large constant $k$, then most sets of maximal twins for $\mathcal{F}$ must have size divisible by $\ell$, provided that the dimension of $\langle \mathcal{F}\rangle_p$ is large. We start with the case when $\ell$ is a prime.
	
	\begin{lemma}\label{lemma:prime}
		Let $V<\mathbb{F}_p^{n}$, let $A_1,\dots,A_d$ be a partition of $[n]$ into twins for $V$, and suppose that $V$ is 2-closed. If $\dim(V) \geq d-h$, then at least $d-2h$ of the numbers $|A_1|,\dots,|A_{d}|$ are divisible by $p$.
	\end{lemma}
	
	\begin{proof}
		For $i\in [d]$, let $b_i=|A_i|$ and let $b$ be the bilinear form defined as in Lemma \ref{lemma:bilinear}. Let $\phi:V\rightarrow\mathbb{F}_p^{d}$ be the linear map defined as $\phi(v)(i)=s$ if $v|_{A_i}$ is the constant $s$ vector ($i = 1,\dots,d$). Then $\phi$ is an injection, so $\dim(\phi(V))=\dim(V)=d-h$. Also, by the definition of $b$, for every $u,v\in V$ we have $||u\cdot v||= \sum_{i = 1}^n{u(i) \cdot v(i)} = \sum_{i=1}^d{|A_i| \cdot \phi(u)(i) \cdot \phi(v)(i)} =  b(\phi(u),\phi(v))$, so we have $b(x,y)=0$ for every $x,y\in\phi(V)$. But then by Lemma \ref{lemma:bilinear}, if $z$ is the number of zeros among $b_1,\dots,b_d$, then $\dim(\phi(V))\leq \frac{1}{2}(d+z)$, which gives $z\geq d-2h$.
	\end{proof}
	
	\begin{lemma}\label{lemma:primepower}
		Let $p$ be a prime and $\alpha\in\mathbb{Z}^{+}$. Let $\mathcal{F}\subset \{0,1\}^{n}$ be $2(p+\alpha)$-closed over $\mathbb{Z}_{p^{\alpha}}$, let $\dim(\langle \mathcal{F}\rangle_p)=d$, and let $A_1,\dots,A_d,B$ be a partition of $[n]$ such that $A_{i}$ is a set of twins for $\mathcal{F}$, and $\dim(\langle \mathcal{F}|_{B}\rangle_{p})\leq h$.  Then at least $d-2\alpha h$ of the numbers $|A_1|,\dots,|A_d|$ are divisible by $p^{\alpha}$.
	\end{lemma}
	
	\begin{proof}
		We prove this by induction on $\alpha$. The case $\alpha=1$ follows from by applying Lemma \ref{lemma:prime} to 
		$V = \langle \mathcal{F}|_{[n] \setminus B}\rangle_{p}$, 
		noting that 
		$\dim(V)\geq \dim(\langle \mathcal{F}\rangle_p) - \dim(\langle \mathcal{F}|_{B}\rangle_{p}) \geq d-h$.
		
		Suppose now that $\alpha>1$. Then, by our induction hypothesis, at least $k=d-(2\alpha-2)h$ of the sets $A_1,\dots,A_d$ have size divisible by $p^{\alpha-1}$, without loss of generality, let these sets be $A_1,\dots,A_k$. Also, let $B'=A_{k+1}\cup\dots\cup A_{d}\cup B$. Letting $V=\langle \mathcal{F}\rangle_p$, note that $$\dim(V|_{B'})\leq \dim(V|_{B})+(d-k)\leq h+d-k.$$ 
		%
		Let $W < V$ be the subspace of all $v \in V$ which vanish on $B'$. Then $\dim(W)\geq d-\dim(V|_{B'})\geq k-h$.
		
		Note that for every $w\in W$ there exists some $w'\in \langle\mathcal{F}\rangle_{p^{\alpha}}$ such that $w'\equiv w \pmod p$. (Indeed, $w \in \langle\mathcal{F}\rangle_{p}$, so $w$ is a linear combination of elements of $\mathcal{F}$ with coefficients in $\mathbb{F}_p$. Taking the same linear combination but over $\mathbb{Z}_{p^{\alpha}}$ gives $w'$.) Let $\beta$ be the smallest number such that $\beta \geq \alpha$ and $\beta=1 \pmod{p-1}$, then $\beta<\alpha+p$. As $\mathcal{F}$ is $2\beta$-closed over $\mathbb{Z}_{p^{\alpha}}$, for every $u,v\in W$ we have that $||(u')^{\beta}\cdot (v')^{\beta}||$ is divisible by $p^{\alpha}$. Note also that we have the following two properties:
		\begin{enumerate}
		    \item[(a)]$(w')^{\beta}\equiv w \pmod p$, as $\beta \equiv 1 \pmod{p-1}$;
		    \item[(b)]If $w(i)=0$ (over $\mathbb{F}_p$) for some $i\in [n]$, then $p^{\alpha}\mid (w')^{\beta}(i)$, because $\beta \geq \alpha$. This means that $(w')^{\beta}(i) \equiv 0 \pmod{p^{\alpha}}$ for every $i \in B'$, since $w$ vanishes on $B'$.
		\end{enumerate}
		For $i\in [k]$, let $A_{i}'$ be a set of size $|A_{i}|/p^{\alpha-1}$, and let $A'=\bigcup_{i=1}^{k}A_{i}'$. Define the linear map $\phi:W\rightarrow\mathbb{F}_p^{A'}$ as follows. If $w\in W$, $i\in [k]$, and $w|_{A_i}$ is the constant $s$ vector, then $\phi(w)|_{A_i'}$ is the constant $s$ vector. 
		(Every $w \in W$ is constant on each $A_i$ because $A_i$ is a set of twins for $V$, and $w$ is determined by its values on $A_1,\dots,A_k$ as $w$ vanishes on $B' = [n] \setminus (A_1 \cup \dots \cup A_k)$.)
		Then $\phi$ is an injection, so $\dim(W)=\dim(\phi(W)) \geq k - h$. Also, for every $u,v\in W$, we have 
		$$||(u')^{\beta}\cdot (v')^{\beta}||\equiv p^{\alpha-1}||\phi(u)\cdot \phi(v)|| \pmod {p^{\alpha}}.$$
		Here we use properties (a)-(b).
		So we see that $||x\cdot y||=0$ (over $\mathbb{F}_p$) for every $x,y\in\phi(W)$, so $\phi(W)$ is $2$-closed. Let $z$ be the number of sets among $A_1',\dots,A_k'$, whose size is divisible by $p$. We can apply Lemma \ref{lemma:prime} to the space $\phi(W)$ and the sets $A'_1,\dots,A'_k$ (recalling that $\dim(\phi(W)) \geq k - h$) to conclude that 
		%
		$z\geq k-2h\geq d-2\alpha h.$
		As $z$ is also the number of sets among $A_1,\dots,A_k$ whose size is divisible by $p^{\alpha}$, this finishes the proof.
	\end{proof}
	
	\begin{lemma}\label{lemma:smalldim}
		Let $p$ be a prime, and $\alpha,t\in\mathbb{Z}^{+}$. Let $\mathcal{F}\subset \{0,1\}^{n}$ such that $\mathcal{F}$ is non-reducible and ${2^{t+1}(p+\alpha)}$-closed over $\mathbb{Z}_{p^{\alpha}}$. Let $A_1,\dots,A_d$ be the unique partition of $[n]$ into maximal sets of twins, and \nolinebreak let $$B=\bigcup_{\substack{i\in [d]\\|A_i|\not\equiv 0\ \mathrm{(mod }\ p^{\alpha})}}A_{i}.$$ Then $\dim(\langle \mathcal{F}|_{B}\rangle_{p})\leq \frac{6n\alpha}{t}$.
	\end{lemma}
	
	\begin{proof}
		Let $\mathcal{F}_0=\mathcal{F}$, and for $i=1,2,\dots,t$, let $\mathcal{F}_{i}=\mathcal{F}_{i-1}\cdot\mathcal{F}_{i-1}$. Note that $\mathcal{F}_{i-1} \subset \mathcal{F}_i$, as it is clear that $\mathcal{F}_{i-1}$ contains only 0-1 vectors. Moreover, $\mathcal{F}_i$ is $2^{t+1-i}(p+\alpha)$-closed over $\mathbb{Z}_{p^{\alpha}}$, which follows by induction and Item 3 in Claim \ref{claim:k-closed}. From this, we will only use that $\mathcal{F}_i$ is at least $2(p+\alpha)$-closed. Finally, $A_1,\dots,A_d$ is also the unique partition of $[n]$ into maximal sets of twins for $\mathcal{F}_i$, see Item 2 in Claim \ref{claim:twins}.
		
		The sequence of positive integers $(\dim(\langle \mathcal{F}_r\rangle_p))_{r=0,\dots,t}$ is monotone increasing (using that $\mathcal{F}_{0}\subset\dots\subset \mathcal{F}_{t})$, and its elements are contained in $[n]$, so by the pigeonhole principle, there exists $0\leq r<t$ such that 
		$$\dim(\langle\mathcal{F}_{r+1} \rangle_{p})\leq \dim(\langle\mathcal{F}_r\rangle_{p})+\frac{n}{t}.$$
		Let $d'=\dim(\langle\mathcal{F}_r\rangle_{p})$. Applying Theorem \ref{cor:stability}, we deduce that $[n]$ can be partitioned into
%
        $d'+1$ sets, $d'$ of which are maximal sets of twins for $\mathcal{F}_r$, and an additional set $C$ satisfying $\dim(\langle \mathcal{F}_r|_{C}\rangle_p)\leq \frac{2n}{t}$. Each maximal set of twins equals $A_i$ for some $i$, so let us assume without loss of generality that the $d'$ maximal sets of twins are $A_1,\dots,A_{d'}$.
		Now, as $\mathcal{F}_{r}$ is $2(p+\alpha)$-closed, we can apply Lemma \ref{lemma:primepower} with $h =  \frac{2n}{t}$ to conclude that at least $q=d'-\frac{4n\alpha}{t}$ of the numbers $|A_{1}|,\dots,|A_{d'}|$ are divisible by $p^{\alpha}$. Without loss of generality, let $A_1,\dots,A_q$ be the sets of twins whose sizes are divisible by $p^{\alpha}$. Let $D=C\cup A_{q+1}\cup\dots\cup A_{d'}$. Then $B\subset D$, and noting that $\dim(\langle\mathcal{F}_r|_{D\setminus C}\rangle_p)\leq d'-q$, and  $\mathcal{F}\subset \mathcal{F}_r$, we get the chain of inequalities $$\dim(\langle \mathcal{F}|_{B}\rangle_p)\leq \dim(\langle \mathcal{F}_{r}|_{D}\rangle_p)\leq (d'-q)+\dim(\langle \mathcal{F}_r|_{C}\rangle_p)\leq \frac{4n\alpha+2n}{t} \leq \frac{6n\alpha}{t}.$$ 
		This finishes the proof.
	\end{proof}
	
	The final ingredient we need for the proof of Theorem \ref{thm:main} is the following well known result, see e.g. the work of Odlyzko \cite{O81}.
	
	\begin{lemma}\label{lemma:odim}
		Let $p$ be a prime and $V< \mathbb{F}_p^{n}$. Then $|V\cap\{0,1\}^{n}|\leq 2^{\dim(V)}$.
	\end{lemma}
	
	\begin{proof}[Proof of Theorem \ref{thm:main}]
		Write $\ell=p_1^{\alpha_1}\dots p_s^{\alpha_s}$, where $p_1,\dots,p_s$ are distinct primes. We show that choosing $k=2^{t+1}\max_{r\in [s]}(p_r+\alpha_r)$ suffices, where $t=12\ell\sum_{r=1}^{s}\alpha_{r}$. More precisely, we show the following two statements. Let $\mathcal{F}\subset \{0,1\}^{n}$ such that $\mathcal{F}$ is $k$-closed over $\mathbb{Z}_{\ell}$.
		\begin{description}
			\item[(1)] Then $|\mathcal{F}|\leq 2^{\lfloor n/\ell\rfloor}$.
			\item[(2)] If $|\mathcal{F}|>2^{\lfloor n/\ell\rfloor-1}$, then $[n]$ can be partitioned into sets $A_{1},\dots,A_{\lfloor n/\ell\rfloor},A'$ such that $A_{i}$ is a maximal set of twins and $|A_i| = \ell$ for each $1 \leq i \leq \lfloor n/\ell\rfloor$, $|A'|\leq \ell-1$, and $\mathcal{F}$ vanishes on $A'$.
		\end{description}
		
		
		
		We proceed by induction on $n$. If $n\leq 6s\ell$, the statements are easy to show. Indeed, let $A_1,\dots,A_d,A'$ be a partition of $[n]$ such that $A_{i}$ is a maximal set of twins for $\mathcal{F}$ for $i\in [d]$, and $\mathcal{F}$ vanishes on $A'$. Then $|\mathcal{F}|\leq 2^{d}$. We argue that as $k>2^{t}\geq 2^{12\ell s}> 6s\ell\geq n$, the  characteristic vector of each $A_{i}$ is contained in $(\langle \mathcal{F} \rangle_{\ell})^{k}$. So fix any $1 \leq i \leq d$. First of all, note that there exists  $v_i\in\mathcal{F}$ such that $v_{i}|_{A_i}$ is the all 1 vector. Now for $j\in [d]\setminus \{i\}$, define $v_j$ as follows. 
		If $v_{i}|_{A_j}$ is 0, then let $v_j=v_i$. Otherwise, as the elements of $A_i$ are not twins of the elements of $A_j$, there exists $w_{j}\in\mathcal{F}$ such that one of $w_j|_{A_{i}}$ and $w_{j}|_{A_{j}}$ is the all 1 vector, and the other is 0. If $w_{j}|_{A_j}$ is 0, then set $v_j=w_j$, otherwise set $v_j=v_i-w_j$. In all cases, $v_j$ has the property that $v_{j}|_{A_i}$ is the all 1 vector, and $v_{j}|_{A_j}$ is 0. But then $v_1\cdot \dotsc \cdot v_{d}$ is the characteristic vector of $A_i$, and as $d\leq n<k$, we have $v_1\cdot\dotsc\cdot v_d\in (\langle \mathcal{F} \rangle_{\ell})^{k}$. Now, since $\mathcal{F}$ is $k$-closed, this implies that $\ell$ divides $|A_{i}|$ for $i\in [d]$. But then $d\leq \lfloor n/\ell\rfloor$, and we are done with (1). Also, if $|\mathcal{F}|>2^{\lfloor n/\ell\rfloor-1}$, we must have $d=\lfloor n/\ell\rfloor$, which is only possible if all the sets $A_1,\dots,A_d$ have size $\ell$. Therefore, (2) also holds.
		
	    From now on assume that $n>6s\ell$. First, suppose that there exists $A\subset [n]$ such that $\ell$ divides $|A|$ and $A$ is a set of twins for $\mathcal{F}$. Then the family $\mathcal{F}'=\mathcal{F}|_{[n]\setminus A}$ is also $k$-closed over $\mathbb{Z}_{\ell}$ and $|\mathcal{F}'|\geq \frac{1}{2}|\mathcal{F}|$. By our induction hypothesis, we have $|\mathcal{F}'|\leq 2^{\lfloor (n-|A|)/\ell\rfloor}\leq 2^{\lfloor (n-\ell)/\ell\rfloor}$, so we get $|\mathcal{F}|\leq 2^{\lfloor n/\ell\rfloor}$, and (1) indeed holds. If $|\mathcal{F}|> 2^{\lfloor n/\ell\rfloor-1}$, then $|\mathcal{F}'|>2^{\lfloor (n-\ell)/\ell\rfloor-1}$, so by our induction hypothesis there exists a partition of $[n]\setminus A$ into sets $A_{1},\dots,A_{\lfloor (n-\ell)/\ell\rfloor},A'$ satisfying (2) with respect to $\mathcal{F}'$. Setting $A_{\lfloor n/\ell\rfloor}=A$, the sets  $A_{1},\dots,A_{\lfloor n/\ell\rfloor},A'$ satisfy (2) with respect to $\mathcal{F}$.
		
		Therefore, in order to finish the proof, it is enough to show that if $|\mathcal{F}|>2^{\lfloor n/\ell\rfloor-1}$, then $\mathcal{F}$ has a set of twins of size divisible by $\ell$. Next, we show that if $I\subset [n]$ is large, then the dimension of $\langle \mathcal{F}|_{I}\rangle_p$ cannot be too small for any prime $p$.
		
		\begin{claim}\label{claim:coordinates}
			Let $p$ be a prime and $I\subset [n]$ such that $|I|\geq \ell$. Then 
			$$|I|\leq \ell\dim(\langle \mathcal{F}|_I\rangle_{p})+3\ell.$$
		\end{claim}
		
		\begin{proof}
			Let $V=\langle \mathcal{F}|_{I}\rangle_p$ and $d = \dim(V)$. Then $|V\cap \{0,1\}^{I}|\leq 2^{d}$ by Lemma \ref{lemma:odim}. By the pigeonhole principle, there exists some $v\in \{0,1\}^{I}$ and $\mathcal{F}'\subset \mathcal{F}$ such that $w|_{I}=v$ for every $w\in\mathcal{F}'$, and $|\mathcal{F}'|\geq |\mathcal{F}|/2^{d}$. Let $0\leq m\leq \ell-1$ such that $||v||\equiv m \pmod \ell$, and for every $w\in \mathcal{F}'$, let $w^{*}$ be the vector we get after replacing the coordinates in $I$ with $m$ coordinates of $1$ entries. This gives a family $\mathcal{F}''=\{w^{*}:w\in\mathcal{F}'\}\subset \{0,1\}^{n-|I|+m}$ such that $|\mathcal{F}''|=|\mathcal{F}'|\geq |\mathcal{F}|/2^{d}$ and $\mathcal{F}''$ is $k$-closed over $\mathbb{Z}_{\ell}$. The latter is true because if $w_1,\dots,w_k\in\mathcal{F}'$, \nolinebreak then  
			$$||w_1^{*}\cdot\dotsc \cdot w_{k}^{*}||=||(w_1|_{[n]\setminus I})\cdot\dotsc\cdot (w_{k}|_{[n]\setminus I})||+m\equiv||w_1\cdot\dotsc\cdot w_{k}||\equiv 0 \pmod \ell .$$
			Therefore, by our induction hypothesis, we have
			$$2^{\lfloor n/\ell\rfloor-d-1}<\frac{|\mathcal{F}|}{2^{d}}\leq |\mathcal{F}''|\leq 2^{\lfloor (n-|I|+m)/\ell\rfloor}< 2^{\lfloor n/\ell\rfloor+2-|I|/\ell}.$$
			Comparing the left- and right-hand-side gives the desired inequality $|I|\leq \ell d+3\ell$. 
		\end{proof}	
		
		We continue with the proof of the theorem.
		We can assume that $\mathcal{F}$ is non-reducible, otherwise we are immediately done 
		by applying our induction hypothesis. Let $A_{1},\dots,A_{d}$ be the unique partition of $[n]$ such that $A_{i}$ is a maximal set of twins for $\mathcal{F}$. Let $r\in [s]$, and apply Lemma \ref{lemma:smalldim} to $\mathcal{F}$ with respect to the prime power $p_r^{\alpha_r}$. Let $$B_r=\bigcup_{\substack{i\in [d]\\|A_i|\not\equiv 0\ \mathrm{(mod }\ p_r^{\alpha_r})}}A_{i}.$$  As $\mathcal{F}$ is $2^{t+1}(p_r+\alpha_r)$-closed, we get from Lemma \ref{lemma:smalldim} that $\dim(\langle \mathcal{F}|_{B_r}\rangle_{p_r})\leq \frac{6n\alpha_r}{t}$. But then by Claim \ref{claim:coordinates}, we also have 
		$$|B_r|\leq \frac{6n\alpha_r\ell}{t}+3\ell.$$
		Let $B=\bigcup_{r=1}^{s}B_{r}$, then 
		$$|B|\leq\sum_{r=1}^{s}|B_{r}|\leq 3s\ell+\frac{6n\ell}{t}\sum_{r=1}^{s}\alpha_r <n,$$
		where the last inequality holds by the choice of $t$ and noting that $n > 6sl$. Observe that $B$ is the union of those maximal sets of twins $A_{i}$ where $|A_i|$ is not divisible by $\ell$. Therefore, as $|B|< n$ and $A_{1},\dots,A_{d}$ form a partition of $[n]$, there must exist $i\in [d]$ such that $\ell$ divides $|A_{i}|$, finishing the proof.
	\end{proof}
